%%% FIELD proof-paired-factorization
Fix \(P\in\mathcal M\), \(x\in\mathcal X\), and \((s,t)\in\mathbb R^2\).  Membership in \(\mathcal M\), error--latent independence, coordinate independence, and \(S_j=A+U_j\) give
\[
\begin{aligned}
F_x(s,t)
&=E_P\!\left[\mathbf 1\{X=x\}e^{i(s+t)A}
 E_P(e^{isU_1})E_P(e^{itU_2})\right]\\
&=p_x\varphi_{f,x}(s+t)\varphi_1(s)\varphi_2(t).
\end{aligned}
\]
The model-class consistency and bounded-potential-outcome conditions imply \(|Y|\le B_Y\) almost surely, so the marked expectation is integrable.  The same conditioning, retaining the mark, yields
\[
\begin{aligned}
G_x(s,t)
&=E_P\!\left[\mathbf 1\{X=x\}Y e^{i(s+t)A}\right]
  \varphi_1(s)\varphi_2(t)\\
&=p_x\left\{\int_{\mathcal A}e^{i(s+t)a}m_x(a)f_x(a)\,da\right\}
  \varphi_1(s)\varphi_2(t)\\
&=p_x\varphi_{q,x}(s+t)\varphi_1(s)\varphi_2(t),
\end{aligned}
\]
where the second equality uses the definitions of \(m_x\) and \(f_x\), and the last uses \(q_x=m_xf_x\).

For protection, \(|e^{iv}-1|\le |v|\), Cauchy--Schwarz, and the error second-moment condition imply, for \(|u|\le\nu_*=(2\sqrt{M_e})^{-1}\),
\[
\begin{aligned}
|\varphi_j(u)|
&\ge 1-E_P|e^{iuU_j}-1|\\
&\ge 1-|u|E_P|U_j|\\
&\ge 1-|u|\{E_P(U_j^2)\}^{1/2}
\ge \frac12,
\end{aligned}
\qquad j\in\{1,2\}.
\]
This proves both channel factorizations and the protected-frequency lower bound.

%%% FIELD proof-adrf-identification
The observed law gives \(p_x=P_{\mathrm{obs}}(X=x)\).  Fix \(x\in\mathcal X\) and put
\[
\rho=\min\{\nu_*/4,(8B_A)^{-1}\}>0.
\]
For \(|s|,|t|\le2\rho\), compact dose support and \(|e^{iv}-1|\le |v|\) imply
\[
|\varphi_{f,x}(s+t)|
\ge 1-E_P|e^{i(s+t)A}-1|
\ge 1-B_A|s+t|\ge\frac12.
\]
The paired factorization, stratum positivity, and the protected-frequency bound therefore give
\[
|F_x(s,t)|\ge p_{\min}/8,
\qquad |s|,|t|\le2\rho. \tag{1}
\]

The second-moment condition implies \(E_P|U_2|<\infty\), and compact dose support implies \(E_P|A|<\infty\), so differentiation under the expectation is justified by domination.  Differentiating the factorization with respect to \(t\) at \((s,0)\), using \(\varphi_2(0)=1\) and the centering condition \(\varphi_2'(0)=iE_PU_2=0\), gives
\[
\frac{\partial_tF_x(s,0)}{F_x(s,0)}
=\frac{\varphi_{f,x}'(s)}{\varphi_{f,x}(s)},
\qquad |s|\le2\rho. \tag{2}
\]
The scalar differential equation in (2), with \(\varphi_{f,x}(0)=1\), has the unique solution
\[
\varphi_{f,x}(s)=
\exp\!\left\{\int_0^s
\frac{\partial_tF_x(u,0)}{F_x(u,0)}\,du\right\},
\qquad |s|\le2\rho. \tag{3}
\]
Thus the observed law identifies \(\varphi_{f,x}\) on a nondegenerate interval.  Because \(f_x\) is integrable and supported on \(\mathcal A\), differentiation under its defining integral at every complex argument shows that \(\varphi_{f,x}\) is entire.  If two candidate transforms agree in (3), their difference is entire and vanishes on a real interval.  Expanding at an interior point and repeatedly taking difference quotients shows that every Taylor coefficient is zero; hence the transforms agree on \(\mathbb C\).  Finally, because \(f_x\in H^\beta\) and \(\beta>1/2\), Cauchy--Schwarz gives
\[
\int_{\mathbb R}|\varphi_{f,x}(u)|\,du
\le
\left\{\int_{\mathbb R}(1+u^2)^{-\beta}\,du\right\}^{1/2}
\|f_x\|_{H^\beta}<\infty. \tag{4}
\]
Convolution with Gaussian approximate identities and passage to the limit in (4) gives the pointwise Fourier inversion formula.  Consequently the continuous version of \(f_x\) is identified on \(\mathbb R\).

On the same interval, (1) permits cancellation of the common nuisance factor:
\[
\frac{G_x(s,0)}{F_x(s,0)}
=\frac{\varphi_{q,x}(s)}{\varphi_{f,x}(s)}. \tag{5}
\]
Equations (3)--(5) identify \(\varphi_{q,x}\) locally.  The model-class outcome bound gives \(|q_x|\le B_Yf_x\), so \(q_x\) is integrable and supported on \(\mathcal A\).  Its transform is therefore entire.  The same analytic identity argument, followed by the Fourier inversion argument using \(q_x\in H^\beta\), identifies the pointwise continuous version of \(q_x\).

It remains to relate this observed-law functional to the causal target.  Consistency and conditional unconfoundedness give, for almost every dose under the conditional density,
\[
E_P\{Y(a)\mid X=x\}
=E_P(Y\mid A=a,X=x)
=m_x(a)=\frac{q_x(a)}{f_x(a)}. \tag{6}
\]
Division on \(I\) is valid because dose overlap gives \(f_x(a)\ge\kappa>0\).  Both \(f_x\) and \(q_x\) are continuous by (4), so the quotient in (6) is continuous on \(I\); the causal-mean-continuity condition supplies a continuous version of the left side.  Since \(f_x\ge\kappa\), an equality holding for conditional-dose almost every \(a\in I\) holds for Lebesgue almost every \(a\in I\), and equality of the two continuous versions extends it to every \(a\in I\).  Averaging over the finite strata now proves
\[
\mu(a)=E_PY(a)
=\sum_{x\in\mathcal X}p_xE_P\{Y(a)\mid X=x\}
=\sum_{x\in\mathcal X}p_x\frac{q_x(a)}{f_x(a)},
\qquad a\in I.
\]
Every quantity on the right has already been recovered from \(P_{\mathrm{obs}}\), which proves identification without confusing the causal curve with its identifying functional.

%%% FIELD statement-exact-profile-compact-measurability
Endow the image of \(\mathcal M\) under \(Q\mapsto((p_{x,Q},f_{x,Q},q_{x,Q})_{x\in\mathcal X},\mathcal L_Q(U_1),\mathcal L_Q(U_2))\) with the product of the Euclidean topology for the stratum masses, the uniform topology on \(\mathcal A\) for the density coordinates, and the weak topology for the two error laws.  This reduced-coordinate image is compact.  For every realized sample, \(Q\mapsto\mathscr D_n(Q)\) and \((Q,a)\mapsto\mu_Q(a)\) are continuous on that image.  The exact profile \(\Theta_{n,\alpha}\), interpreted through reduced coordinates, is a measurable compact-valued correspondence, possibly empty, and the functions \((O_1,\ldots,O_n,a)\mapsto\inf_{Q\in\Theta_{n,\alpha}}\mu_Q(a)\) and \((O_1,\ldots,O_n,a)\mapsto\sup_{Q\in\Theta_{n,\alpha}}\mu_Q(a)\), with both set equal to zero when the profile is empty, are jointly measurable on the sample space times \(I\).

%%% FIELD proof-exact-profile-compact-measurability
We first prove compactness rather than assume it.  Consider any sequence of reduced coordinates arising from laws in \(\mathcal M\).  The stratum-mass vectors lie in the closed simplex cut out by \(p_x\ge p_{\min}\), hence have a convergent subsequence.  For any function \(h\) supported on \(\mathcal A\) with \(\|h\|_{H^\beta}\le R\), Fourier inversion and Cauchy--Schwarz give the common bound
\[
\|h\|_\infty
\le C R\left\{\int_{\mathbb R}(1+t^2)^{-\beta}\,dt\right\}^{1/2}<\infty. \tag{7}
\]
Moreover, for \(a,a'\in\mathcal A\),
\[
|h(a)-h(a')|
\le CR\left[\int_{\mathbb R}
\min\{4,t^2|a-a'|^2\}(1+t^2)^{-\beta}\,dt\right]^{1/2}. \tag{8}
\]
The integral in (8) tends to zero with \(|a-a'|\) by dominated convergence, since \(\beta>1/2\).  Thus both the density and marked-density balls are uniformly bounded and equicontinuous on \(\mathcal A\).  A diagonal subsequence on a countable dense subset, followed by (8), converges uniformly.  Boundedness in the Hilbert space \(H^\beta\) also gives a weakly convergent subsubsequence.  Its weak limit equals the uniform limit as a distribution, and weak lower semicontinuity preserves the respective \(H^\beta\) bounds.  Nonnegativity, unit integral, support in \(\mathcal A\), overlap on \(I\), and the pointwise relation \(|q_x|\le B_Yf_x\) are closed under uniform convergence.

For either error coordinate, Markov's inequality gives uniformly
\[
Q(|U_j|>T)\le M_eT^{-2}. \tag{9}
\]
Choosing successively convergent distribution functions on rational continuity points and using (9) to prevent loss of mass produces a weakly convergent subsequence.  The function \(u\mapsto u^2\) is the increasing limit of bounded continuous truncations, so the limit law has second moment at most \(M_e\).  In addition,
\[
E_Q\{|U_j|\mathbf 1(|U_j|>T)\}\le M_e/T, \tag{10}
\]
so first moments are uniformly integrable.  Hence the weak limit of the second error remains centered.  This proves sequential compactness of all coordinate factors; they are metrizable, so it proves compactness once closedness of the image is checked.

For that check, let a convergent coordinate sequence have limit \((p_x,f_x,q_x,\nu_1,\nu_2)\).  Put \(m_x=q_x/f_x\) where \(f_x>0\), and put \(m_x=0\) where \(f_x=0\).  The preserved inequality \(|q_x|\le B_Yf_x\) gives \(|m_x|\le B_Y\), and overlap makes \(m_x\) continuous on \(I\).  Construct \(X\) with masses \(p_x\), draw \(A\mid X=x\) with density \(f_x\), draw independent errors with laws \(\nu_1,\nu_2\), independently of \((A,X)\), and define the potential-outcome schedule deterministically by \(Y(a)=m_X(a)\), with \(Y=Y(A)\).  This construction satisfies consistency, conditional unconfoundedness, bounded outcomes, causal-mean continuity on \(I\), both error-independence conditions, centering, all support and Sobolev restrictions, and both overlap restrictions.  Its marked density is \(q_x\).  Thus the limiting coordinates are again the image of a member of \(\mathcal M\), proving compactness of the reduced image.

We next prove continuity.  Weak convergence of error laws implies pointwise convergence of their characteristic functions.  The common bound \(E|U_j|\le\sqrt{M_e}\) makes these characteristic functions uniformly Lipschitz, so a finite-net argument upgrades pointwise convergence to uniform convergence on every compact frequency interval.  Uniform convergence of \(f_x\) and \(q_x\) on \(\mathcal A\) gives uniform convergence of their Fourier transforms on compact frequency intervals.  The paired channel factorization therefore gives uniform convergence of \(F_{x,Q}\) and \(G_{x,Q}\) on \(\mathcal C_*\).  Integration over the finite-volume cube and convergence of the mass coordinates prove continuity of \(Q\mapsto\mathscr D_n(Q)\) for each fixed sample.  Finally, overlap gives, uniformly on \(I\),
\[
\left|\frac{q_x}{f_x}-\frac{\widetilde q_x}{\widetilde f_x}\right|
\le \kappa^{-1}|q_x-\widetilde q_x|
+B_Y\kappa^{-1}|f_x-\widetilde f_x|, \tag{11}
\]
so the finite identifying sum proves joint continuity of \((Q,a)\mapsto\mu_Q(a)\).

Let the compact reduced image just proved be denoted temporarily by \(K\) inside this proof.  The map \((O_1,\ldots,O_n,z)\mapsto\mathscr D_n(z)\) is measurable in the sample and continuous in \(z\): this follows directly from the finite empirical sums and the preceding continuity argument.  Hence the sublevel set \(\Theta_{n,\alpha}=\{z\in K:\mathscr D_n(z)\le\epsilon_n^2\}\) is compact for every sample.  To verify correspondence measurability directly, let \(V\) be open in \(K\), and define the compact sets
\[
C_m(V)=\{z\in K:\operatorname{dist}(z,K\setminus V)\ge m^{-1}\}.
\]
Then \(\Theta_{n,\alpha}\cap V\ne\varnothing\) exactly when, for some \(m\),
\[
\inf_{z\in C_m(V)}\mathscr D_n(z)\le\epsilon_n^2. \tag{12}
\]
For a fixed countable dense subset of each \(C_m(V)\), continuity makes the infimum in (12) the infimum over that countable set, hence measurable.  A countable base for the compact metric space \(K\) completes the correspondence-measurability claim.

It remains to establish joint endpoint measurability without invoking an unstated selection theorem.  For positive integers \(m\), define
\[
v_m^-(o,a)=\inf_{z\in K}\left[\mu_z(a)+m\{\mathscr D_n(o,z)-\epsilon_n^2\}_+\right],
\]
\[
v_m^+(o,a)=\sup_{z\in K}\left[\mu_z(a)-m\{\mathscr D_n(o,z)-\epsilon_n^2\}_+\right]. \tag{13}
\]
Each extremum in (13) equals the corresponding extremum over one fixed countable dense subset of \(K\), because the bracketed objectives are continuous in \(z\).  Thus \(v_m^-\) and \(v_m^+\) are jointly measurable in \((o,a)\).  If the profile is nonempty, compactness and a convergent-subsequence argument for optimizing points show
\[
v_m^-(o,a)\uparrow\inf_{z\in\Theta_{n,\alpha}(o)}\mu_z(a),
\qquad
v_m^+(o,a)\downarrow\sup_{z\in\Theta_{n,\alpha}(o)}\mu_z(a). \tag{14}
\]
Indeed, boundedness of \(\mu_z(a)\) forces every subsequential optimizer with bounded penalized value to have \(\mathscr D_n(o,z)\le\epsilon_n^2\), and continuity then gives the two limiting extrema.  If the profile is empty, compactness makes \(\min_{z\in K}\mathscr D_n(o,z)>\epsilon_n^2\), so this case is itself measurable, and the stipulated zero value is measurable.  The pointwise limits in (14), combined with that empty-profile convention, prove joint measurability of both exact endpoint functions.

%%% FIELD construction-profiled-band-exact
If \(\Theta_{n,\alpha}\neq\varnothing\), set
\[
L_n(a)=\max\left\{-B_Y,\inf_{Q\in\Theta_{n,\alpha}}\mu_Q(a)\right\},
\qquad
U_n(a)=\min\left\{B_Y,\sup_{Q\in\Theta_{n,\alpha}}\mu_Q(a)\right\},
\qquad a\in I.
\]
If \(\Theta_{n,\alpha}=\varnothing\), set \(L_n(a)=U_n(a)=0\) for every \(a\in I\).

%%% FIELD construction-upper-handle-exact
Use the continuous integrated probability and signed marked contrasts on \(\mathcal C_*\), derive their nearly parametric forward radius, pass both compactly supported entire transforms through the quantitative analytic-continuation argument at order \(K_n\), recover the protected quotient \(q_x/f_x\) on \(I\), and take measurable extrema over the exact compact profile \(\Theta_{n,\alpha}\).

%%% FIELD statement-finite-profile-certification-oeq
Under exactly the model class \(\mathcal M\), can one construct from the sample and fixed class constants a terminating algorithm that returns a finite, finitely checkable witness cover of the exact profile \(\Theta_{n,\alpha}\), together with rigorous endpoint enclosures whose deterministic error is \(o(r_n)\), without deciding nonemptiness of an abstract infinite-dimensional decoder?  Such an effective certification result is not established here and remains an open computational question; the statistical band and its rate use the exact profile directly.

%%% FIELD proof-honest-upper-band
For \(P\in\mathcal M\), let
\[
E_P=\{\mathscr D_n(P)\le\epsilon_n^2\}.
\]
The established continuous-channel concentration proposition, under the i.i.d. sampling condition, gives
\[
\inf_{P\in\mathcal M}P_{\mathrm{obs}}^{\otimes n}(E_P)\ge1-\alpha \tag{15}
\]
for all sufficiently large \(n\).  On \(E_P\), the true law belongs to the exact profile \(\Theta_{n,\alpha}\), so that profile is nonempty.  The model-class outcome and consistency conditions give \(|\mu_Q(a)|\le B_Y\) for every \(Q\in\mathcal M\) and \(a\in I\).  Therefore the clipping in the exact-profile definition cannot remove \(\mu_P\), and
\[
L_n(a)\le\mu_P(a)\le U_n(a)
\quad\text{for every }a\in I \quad\text{on }E_P. \tag{16}
\]
The exact-profile compactness and measurability lemma makes \(L_n\) and \(U_n\) jointly measurable.  If the profile is nonempty, its infimum is no larger than its supremum and both lie in \([-B_Y,B_Y]\); if it is empty, both endpoints are zero.  Thus \(L_n\le U_n\).  Equations (15)--(16) prove \((L_n,U_n)\in\mathfrak B_{n,\alpha}\).

For width, view \(\mathscr D_n(Q)^{1/2}\) as the norm of the difference between the empirical channel-and-mass vector and the corresponding vector under \(Q\) in the Hilbert direct sum specified by the definition of \(\mathscr D_n\).  If \(Q,R\in\Theta_{n,\alpha}\), the Hilbert-space triangle inequality gives
\[
\begin{aligned}
&\left[\sum_{x\in\mathcal X}\left\{
\int_{\mathcal C_*}\left(|F_{x,Q}-F_{x,R}|^2
+B_Y^{-2}|G_{x,Q}-G_{x,R}|^2\right)\,ds\,dt
+|p_{x,Q}-p_{x,R}|^2\right\}\right]^{1/2}\\
&\hspace{35mm}\le \mathscr D_n(Q)^{1/2}+\mathscr D_n(R)^{1/2}
\le2\epsilon_n. 
\end{aligned} \tag{17}
\]
The paired protected inverse theorem, applied with \(\delta=2\epsilon_n\), therefore yields
\[
\sup_{Q,R\in\Theta_{n,\alpha}}
\|\mu_Q-\mu_R\|_{L^\infty(I)}
\le C\left\{\frac{\log\log(1/\epsilon_n)}{\log(1/\epsilon_n)}\right\}^{\beta-1/2}. \tag{18}
\]
Since \(\epsilon_n=C_\alpha\sqrt{\log(en/\alpha)/n}\), fixed \(\alpha\) implies
\[
\log(1/\epsilon_n)\asymp\log n,
\qquad
\log\log(1/\epsilon_n)\asymp\log\log n. \tag{19}
\]
The definition of \(K_n\), with fixed \(c_K>0\), then turns (18)--(19) into
\[
\sup_{Q,R\in\Theta_{n,\alpha}}
\|\mu_Q-\mu_R\|_{L^\infty(I)}\le C K_n^{1/2-\beta}=Cr_n. \tag{20}
\]
For a nonempty compact profile, the endpoint lemma makes the extrema attainable; even without choosing optimizers, (20) directly bounds every pointwise diameter and hence gives \(W_n\le Cr_n\).  For an empty profile, \(W_n=0\).  Thus the bound holds for every realized sample, and consequently
\[
\sup_{P\in\mathcal M}E_{P_{\mathrm{obs}}^{\otimes n}}[W_n]\le Cr_n.
\]
This establishes the theoretical exact-profile upper band without any computational certification premise.

%%% FIELD proof-fixed-nuisance-marked-pair
Apply the fixed-kernel Bernoulli packet lemma with
\[
k_n=\left\lceil\frac{\log(en)}{\log\log(e^e n)}\right\rceil
\]
and set \(P_{j,n}=P_{j,k_n}\).  That lemma constructs members of \(\mathcal M\), with common fixed stratum masses, common density \(f^\circ\), common fixed nonidentical asymmetric error laws, and conditionally Bernoulli marks.  Since \(c_K>0\) is fixed,
\[
k_n\asymp K_n,
\qquad
\|\mu_{1,n}-\mu_{0,n}\|_{L^\infty(I)}
\ge c k_n^{1/2-\beta}\ge c r_n. \tag{21}
\]

Write \(\chi_n^2=\chi^2(P_{1,n,\mathrm{obs}},P_{0,n,\mathrm{obs}})\).  Direct multiplication of likelihood ratios gives the product identity
\[
1+\chi^2(P_{1,n,\mathrm{obs}}^{\otimes n},P_{0,n,\mathrm{obs}}^{\otimes n})
=(1+\chi_n^2)^n. \tag{22}
\]
The packet lemma and \(m!\ge(m/e)^m\) imply
\[
n\chi_n^2
\le Cn\frac{C_0^{4k_n}}{\{(2k_n-4)!\}^2}\longrightarrow0, \tag{23}
\]
because the logarithm of the right side is at most
\[
\log n-4k_n\log k_n+O(k_n)\longrightarrow-\infty.
\]
Using \((1+x)^n\le e^{nx}\), equations (22)--(23) show that the product chi-square divergence tends to zero.  For any likelihood ratio \(Z\), Cauchy--Schwarz gives
\[
\|P-Q\|_{\mathrm{TV}}
=\frac12E_Q|Z-1|
\le\frac12\{E_Q(Z-1)^2\}^{1/2},
\]
and hence
\[
\|P_{1,n,\mathrm{obs}}^{\otimes n}-P_{0,n,\mathrm{obs}}^{\otimes n}\|_{\mathrm{TV}}
\le\frac12\left\{\chi^2(P_{1,n,\mathrm{obs}}^{\otimes n},P_{0,n,\mathrm{obs}}^{\otimes n})\right\}^{1/2}
\longrightarrow0. \tag{24}
\]
It is therefore at most the prescribed \(\eta>0\) for all sufficiently large \(n\).  Combining (21) and (24) proves the theorem.

%%% FIELD proof-minimax-width-frontier
The exact-profile upper-band theorem supplies one data-dependent element of \(\mathfrak B_{n,\alpha}\).  Therefore the definition of minimax expected maximal width gives
\[
\mathfrak R_n
\le\sup_{P\in\mathcal M}E_{P_{\mathrm{obs}}^{\otimes n}}[W_n]
\le Cr_n. \tag{25}
\]

For the converse, take an arbitrary \((L_n,U_n)\in\mathfrak B_{n,\alpha}\) and the two laws from the fixed-nuisance marked-pair theorem.  Define
\[
A_j=\{L_n(a)\le\mu_{j,n}(a)\le U_n(a)
\text{ for every }a\in I\},
\qquad j\in\{0,1\}.
\]
Uniform honesty gives
\[
P_{j,n,\mathrm{obs}}^{\otimes n}(A_j)\ge1-\alpha. \tag{26}
\]
The defining variational inequality for total variation gives
\[
P_{0,n,\mathrm{obs}}^{\otimes n}(A_1)
\ge P_{1,n,\mathrm{obs}}^{\otimes n}(A_1)
-\|P_{1,n,\mathrm{obs}}^{\otimes n}-P_{0,n,\mathrm{obs}}^{\otimes n}\|_{\mathrm{TV}}
\ge1-\alpha-\eta. \tag{27}
\]
The union bound and (26)--(27) yield
\[
P_{0,n,\mathrm{obs}}^{\otimes n}(A_0\cap A_1)
\ge1-2\alpha-\eta>0, \tag{28}
\]
where the strict positivity is exactly the declared condition \(\eta<1-2\alpha\).  On \(A_0\cap A_1\), both causal curves lie pointwise in the same band, so
\[
W_n\ge\|\mu_{1,n}-\mu_{0,n}\|_{L^\infty(I)}\ge cr_n. \tag{29}
\]
Taking expectation under \(P_{0,n,\mathrm{obs}}^{\otimes n}\), then the supremum over \(P\in\mathcal M\), and finally the infimum over all honest bands proves
\[
\mathfrak R_n\ge c(1-2\alpha-\eta)r_n. \tag{30}
\]
For fixed \(c_K>0\) and \(\beta>1/2\),
\[
K_n\asymp\frac{\log n}{\log\log n},
\qquad
r_n=K_n^{1/2-\beta}
\asymp\left\{\frac{\log\log n}{\log n}\right\}^{\beta-1/2}. \tag{31}
\]
Equations (25), (30), and (31) prove the matching frontier.

%%% FIELD proof-error-free-reduction
On the stated submodel, \(U_1=U_2=0\) almost surely, so \(\varphi_1(s)=\varphi_2(t)=1\) for every \(s,t\).  Substitution into the paired-channel factorization gives
\[
F_x(s,t)=p_x\varphi_{f,x}(s+t),
\qquad
G_x(s,t)=p_x\varphi_{q,x}(s+t).
\]
The identification theorem remains applicable and, using \(q_x=m_xf_x\) and overlap on \(I\), yields
\[
\mu(a)=\sum_{x\in\mathcal X}p_x\frac{q_x(a)}{f_x(a)}
=\sum_{x\in\mathcal X}p_xE_P[Y\mid A=a,X=x].
\]
Consistency and conditional unconfoundedness identify the last conditional regression with \(E_P\{Y(a)\mid X=x\}\) for almost every \(a\); causal-mean continuity and the continuous quotient extend the equality to every \(a\in I\).  Averaging over \(X\) gives \(E_PY(a)\), which is exactly the standard continuous-treatment g-formula asserted in the theorem.

%%% FIELD proof-paired-protected-inverse
Let \(d(P,Q)\) be the square root of the displayed paired-channel discrepancy and suppose \(d(P,Q)\le\delta\).  Put
\[
\rho=\min\{\nu_*/4,(8B_A)^{-1}\}.
\]
Then \([-2\rho,2\rho]^2\subset\mathcal C_*\).  Under either law, the paired factorization, stratum positivity, compact dose support, and the protected-frequency bound give
\[
|F_x(s,t)|
\ge p_{\min}\left(1-B_A|s+t|\right)\frac12\frac12
\ge p_{\min}/8,
\qquad |s|,|t|\le2\rho. \tag{32}
\]

We first turn integrated channel control into uniform local control.  Every probability channel is uniformly Lipschitz because
\[
E|S_j|\le B_A+\sqrt{M_e},
\]
and every marked channel has the same Lipschitz bound multiplied by \(B_Y\).  If \(H\) is \(L\)-Lipschitz on the outer square and its maximum modulus on the inner square is \(M\), then on the intersection of the outer square with a disk of radius \(M/(2L)\) about a maximizer, \(|H|\ge M/2\).  The inner square has a fixed positive margin, so for sufficiently small \(\|H\|_2\), integration on that disk gives
\[
\|H\|_2^2\ge cM^4/L^2,
\qquad M\le C\|H\|_2^{1/2}. \tag{33}
\]
Applying (33) to each channel difference and using finiteness of \(\mathcal X\) yields
\[
\max_{x\in\mathcal X}\left(
\|F_{x,P}-F_{x,Q}\|_{\infty,[-\rho,\rho]^2}
+B_Y^{-1}\|G_{x,P}-G_{x,Q}\|_{\infty,[-\rho,\rho]^2}
\right)\le C\delta^{1/2}. \tag{34}
\]

The second \(t\)-derivatives of the probability channels are uniformly bounded, since
\[
E S_2^2\le2(B_A^2+M_e).
\]
For a twice differentiable scalar function \(h\), Taylor's formula at \(v\) and \(-v\) gives
\[
|h'(0)|\le \frac{\|h\|_\infty}{v}+\frac v2\|h''\|_\infty.
\]
Choosing \(v\) proportional to \(\|h\|_\infty^{1/2}\), within the fixed outer-square margin, and applying (34) to \(h(t)=F_{x,P}(s,t)-F_{x,Q}(s,t)\), gives
\[
\sup_{|s|\le\rho}
|\partial_tF_{x,P}(s,0)-\partial_tF_{x,Q}(s,0)|
\le C\delta^{1/4}. \tag{35}
\]

Centering of the second error and factorization imply, under each law,
\[
\frac{\partial_tF_x(s,0)}{F_x(s,0)}
=\frac{\varphi_{f,x}'(s)}{\varphi_{f,x}(s)}.
\]
The lower bound (32), together with (34)--(35), bounds the difference of the two logarithmic derivatives by \(C\delta^{1/4}\).  Integrating from zero to \(s\), using \(\varphi_{f,x,P}(0)=\varphi_{f,x,Q}(0)=1\), and using the Lipschitz property of the exponential map on the resulting bounded logarithms yields
\[
\sup_{|s|\le\rho}
|\varphi_{f,x,P}(s)-\varphi_{f,x,Q}(s)|
\le C\delta^{1/4}. \tag{36}
\]
Also, within each law the common nuisance factor cancels:
\[
\frac{G_x(s,0)}{F_x(s,0)}
=\frac{\varphi_{q,x}(s)}{\varphi_{f,x}(s)}.
\]
Using (32), (34), (36), and \(|\varphi_{q,x}(s)|\le\int|q_x|\le B_Y\) gives
\[
\sup_{|s|\le\rho}
|\varphi_{q,x,P}(s)-\varphi_{q,x,Q}(s)|
\le C\delta^{1/4}. \tag{37}
\]

Apply the continuous analytic Sobolev modulus lemma to \(f_{x,P}-f_{x,Q}\) and \(q_{x,P}-q_{x,Q}\).  Their supports lie in \(\mathcal A\), their Sobolev norms are at most \(2R_f\) and \(2R_q\), and their \(L^1\) norms are at most \(2\) and \(2B_Y\), respectively.  Since replacing \(\delta\) by a fixed multiple of \(\delta^{1/4}\) changes the following logarithmic expression only by fixed constants, (36)--(37) imply
\[
\max_{x\in\mathcal X}\left(
\|f_{x,P}-f_{x,Q}\|_\infty+
\|q_{x,P}-q_{x,Q}\|_\infty\right)
\le C\omega(\delta), \tag{38}
\]
where
\[
\omega(\delta)=
\left\{\frac{\log\log(1/\delta)}{\log(1/\delta)}\right\}^{\beta-1/2}.
\]
The discrepancy also directly gives \(\max_x|p_{x,P}-p_{x,Q}|\le\delta\).  Dose overlap and \(|q_{x,Q}|\le B_Yf_{x,Q}\) imply on \(I\)
\[
\left|\frac{q_{x,P}}{f_{x,P}}-
\frac{q_{x,Q}}{f_{x,Q}}\right|
\le \kappa^{-1}|q_{x,P}-q_{x,Q}|
+B_Y\kappa^{-1}|f_{x,P}-f_{x,Q}|. \tag{39}
\]
Combining (38)--(39), the finite identifying sum, and the mass difference proves
\[
\|\mu_P-\mu_Q\|_{L^\infty(I)}\le C\omega(\delta).
\]

For sharpness, apply the fixed-kernel Bernoulli packet lemma with
\[
k(\delta)=\left\lceil c_0
\frac{\log(1/\delta)}{\log\log(1/\delta)}\right\rceil,
\]
where \(c_0\) is a sufficiently large fixed constant.  Since \(m!\ge(m/e)^m\),
\[
C\frac{C_0^{4k(\delta)}}{\{(2k(\delta)-4)!\}^2}\le c_1\delta^2 \tag{40}
\]
for all sufficiently small \(\delta\), with \(c_1\) made as small as needed by increasing \(c_0\).  The packet pair has identical probability channels and identical masses.  If \(H\) is a complex random variable with \(|H|\le1\), the variational definition of total variation gives \(|E_PH-E_QH|\le2\|P-Q\|_{\mathrm{TV}}\); Cauchy--Schwarz gives \(4\|P-Q\|_{\mathrm{TV}}^2\le\chi^2(P,Q)\).  Apply these facts to \(Y e^{i(sS_1+tS_2)}/B_Y\).  Because \(\mathcal C_*\) has finite volume and \(\mathcal X\) is finite, (40) makes the full paired-channel discrepancy at most \(\delta^2\).  The packet separation is
\[
ck(\delta)^{1/2-\beta}\ge c\omega(\delta). \tag{41}
\]
All stratum masses, latent densities, and error laws are fixed across this pair, and its marks are conditionally Bernoulli.  Equations (40)--(41) prove the asserted converse and unimprovability.

%%% FIELD statement-protected-modulus-characterization
There are constants \(0<c<C<\infty\) such that, for every sufficiently small \(\delta>0\),
\[
\begin{aligned}
c\left\{\frac{\log\log(1/\delta)}{\log(1/\delta)}\right\}^{\beta-1/2}
&\le
\sup\Bigg\{\|\mu_P-\mu_Q\|_{L^\infty(I)}:P,Q\in\mathcal M,\\
&\quad \sum_{x\in\mathcal X}\left[\int_{\mathcal C_*}
\left(|F_{x,P}-F_{x,Q}|^2+B_Y^{-2}|G_{x,P}-G_{x,Q}|^2\right)\,ds\,dt
+|p_{x,P}-p_{x,Q}|^2\right]\le\delta^2\Bigg\}\\
&\le C\left\{\frac{\log\log(1/\delta)}{\log(1/\delta)}\right\}^{\beta-1/2}.
\end{aligned}
\]
The lower inequality remains true when the supremum is restricted to pairs with fixed common \(p_x\), fixed common \(f_x=f^\circ\), fixed common nonidentical asymmetric error laws, and conditionally Bernoulli outcomes.

%%% FIELD proof-protected-modulus-characterization
The upper inequality is the first conclusion of the paired protected inverse theorem, applied to every admissible pair and then followed by a supremum.  The converse clause of that same theorem supplies, for every sufficiently small \(\delta\), an admissible pair in the stated fixed-nuisance Bernoulli submodel with separation at least a fixed positive multiple of \(\{\log\log(1/\delta)/\log(1/\delta)\}^{\beta-1/2}\).  Taking the restricted supremum over a set containing that pair proves the lower inequality.  Hence the question has the positive sharp characterization stated here.

%%% FIELD proof-continuous-analytic-sobolev-modulus
Put
\[
H(z)=\int_{-B_A}^{B_A}e^{iza}h(a)\,da.
\]
Absolute convergence on compact subsets of \(\mathbb C\) justifies differentiation under the integral and gives
\[
H^{(j)}(z)=\int_{-B_A}^{B_A}(ia)^je^{iza}h(a)\,da,
\qquad |H^{(j)}(0)|\le MB_A^j. \tag{42}
\]
Let \(T_k(z)=\sum_{j=0}^kH^{(j)}(0)z^j/j!\).  Taylor's formula inside the integral gives
\[
|H(t)-T_k(t)|
\le M\frac{(B_A|t|)^{k+1}}{(k+1)!},
\qquad t\in\mathbb R. \tag{43}
\]

We need one elementary polynomial extrapolation bound.  If \(p\) has degree at most \(k\), rescale \([-\rho,\rho]\) to \([-1,1]\) and expand in the orthonormal Legendre basis.  Cauchy--Schwarz bounds every coefficient by \(\|p\|_{L^2[-\rho,\rho]}\), while the three-term recursion, by induction, gives \(|P_j(x)|\le(2|x|+1)^j\) for \(|x|\ge1\).  Summing over \(0\le j\le k\) yields
\[
\sup_{|t|\le T}|p(t)|
\le C\rho^{-1/2}
\left\{\frac{C(1+T)}{\rho}\right\}^k
\|p\|_{L^2[-\rho,\rho]},
\qquad T\ge\rho. \tag{44}
\]
On \([-\rho,\rho]\), (43) and the hypothesis imply
\[
\|T_k\|_2
\le\varepsilon+(2\rho)^{1/2}M
\frac{(B_A\rho)^{k+1}}{(k+1)!}. \tag{45}
\]
Choose a fixed \(a_0>0\), depending only on \(B_A\) and \(\rho\), sufficiently small and set \(T=a_0k\).  The elementary factorial bound \((k+1)!\ge((k+1)/e)^{k+1}\) shows that both
\[
\left\{\frac{C(1+T)}{\rho}\right\}^k
\frac{(B_A\rho)^{k+1}}{(k+1)!}
\quad\text{and}\quad
\frac{(B_AT)^{k+1}}{(k+1)!}
\]
are at most \(Ce^{-ck}\).  Combining (43)--(45) therefore gives
\[
\sup_{|t|\le T}|H(t)|
\le (C_0k)^k\varepsilon+CMe^{-ck}. \tag{46}
\]
Because \(\beta>1/2\), Cauchy--Schwarz in Fourier space gives
\[
\int_{|t|>T}|H(t)|\,dt
\le
\left\{\int_{|t|>T}(1+t^2)^{-\beta}\,dt\right\}^{1/2}
\|h\|_{H^\beta}
\le CRT^{1/2-\beta}. \tag{47}
\]
Fourier inversion is valid because the right side of (47), together with (46), shows \(H\in L^1\).  Splitting the inversion integral at \(T=a_0k\), bounding its inner part by \(2T\) times (46), and using (47) proves
\[
\|h\|_\infty
\le C\left[Rk^{1/2-\beta}
+k(C_0k)^k\varepsilon+Mk e^{-ck}\right]. \tag{48}
\]

For the consequence, let \(L=\log(1/\varepsilon)\) and choose
\[
k=\left\lfloor c_1\frac{L}{\log(e+L)}\right\rfloor
\]
with \(c_1>0\) sufficiently small.  Then
\[
k(C_0k)^k\varepsilon+Mk e^{-ck}
\le C(R+M)k^{1/2-\beta}
\]
for all sufficiently large \(L\): after taking logarithms, the first term has leading exponent \(-(1-c_1+o(1))L\), and the second decays exponentially in \(k\), whereas \(k^{1/2-\beta}\) is polynomial.  Since \(k\asymp L/\log L\), substitution in (48) gives
\[
\|h\|_\infty
\le C(R+M)
\left\{\frac{\log\log(1/\varepsilon)}{\log(1/\varepsilon)}\right\}^{\beta-1/2}.
\]

%%% FIELD proof-fixed-kernel-bernoulli-packet
Choose a closed nondegenerate interval \(J\subset\operatorname{int}(I)\), map it affinely to \([-1,1]\) by \(a\mapsto z(a)\), fix an integer \(r>\beta+4\), and put \(w(z)=(1-z^2)^r\).  Let \(\{\pi_j:j\ge0\}\) be the orthonormal Jacobi polynomials in \(L^2(w(z)\,dz)\), and define
\[
K_k(z,0)=\sum_{j=2k}^{4k}\pi_j(0)\pi_j(z),
\qquad
\widetilde h_k(a)=w(z(a))K_k(z(a),0)\mathbf 1\{a\in J\}. \tag{49}
\]
Every polynomial of degree below \(2k\) is in the span of \(\pi_0,\ldots,\pi_{2k-1}\).  Orthogonality in (49) therefore gives
\[
\int_J a^\ell\widetilde h_k(a)\,da=0,
\qquad 0\le\ell<2k. \tag{50}
\]

We record the estimates needed below and derive their orders from the Jacobi formula.  Rodrigues' formula and the exact squared norm give, for even degrees,
\[
P_{2m}^{(r,r)}(0)
=(-1)^m\frac{\Gamma(2m+r+1)}{2^{2m}m!\,\Gamma(m+r+1)},
\]
and the ratio of consecutive gamma factors shows that the normalized values \(|\pi_{2m}(0)|\) are bounded above and below by positive constants; odd values vanish.  Hence
\[
c k\le K_k(0,0)\le Ck. \tag{51}
\]
Repeated integration by parts in Rodrigues' formula gives, for every fixed integer \(m\le r-1\),
\[
\|\partial_z^m\{wK_k(\cdot,0)\}\|_2^2
\le C_m\sum_{j=2k}^{4k}(1+j)^{2m}|\pi_j(0)|^2
\le C_m k^{2m+1}. \tag{52}
\]
The same factorial-product bounds in the Christoffel--Darboux identity give
\[
\sup_{z\in[-1,1]}|w(z)K_k(z,0)|\le Ck. \tag{53}
\]
Because \(w\) has a zero of order \(r\) at each endpoint, extension by zero introduces no boundary distribution through derivative order \(r-1\).  After the fixed affine change of variables, (51)--(53) imply
\[
\|\widetilde h_k\|_2\le Ck^{1/2},
\qquad
\|\widetilde h_k\|_{H^m}\le C_mk^{m+1/2},
\qquad
\|\widetilde h_k\|_\infty\le Ck. \tag{54}
\]
If \(m>\beta\), Hölder's inequality applied to
\((1+t^2)^\beta|\widehat{\widetilde h_k}(t)|^2\) gives the interpolation inequality
\[
\|\widetilde h_k\|_{H^\beta}
\le\|\widetilde h_k\|_2^{1-\beta/m}
\|\widetilde h_k\|_{H^m}^{\beta/m}
\le Ck^{\beta+1/2}. \tag{55}
\]
After multiplication by one fixed sufficiently small constant, set
\[
h_k=k^{-\beta-1/2}\widetilde h_k.
\]
Equations (51), (54), and (55), together with Cauchy--Schwarz on the fixed interval \(J\), yield
\[
\|h_k\|_{H^\beta}\le1,
\quad \|h_k\|_1\le Ck^{-\beta},
\quad \|h_k\|_\infty\le Ck^{1/2-\beta},
\quad |h_k(a_0)|\ge ck^{1/2-\beta}, \tag{56}
\]
where \(a_0\) corresponds to \(z=0\).  The moment cancellations (50) are unchanged.

We next construct two fixed error laws.  Define
\[
\psi(u)=\left\{\frac{\sin(u/2)}{u/2}\right\}^4,
\qquad
g_{a,v}^{\sharp}(u)=c_{a,v}\{\psi(u)+v\psi(u-a)\}, \tag{57}
\]
where \(a\notin2\pi\mathbb Z\), \(v>0\), \(v\ne1\), and \(c_{a,v}\) normalizes the integral.  This is a nonnegative asymmetric density.  The Fourier transform of \(\psi\) is the fourfold convolution of compact interval indicators, so the characteristic function of (57) has compact support.  The functions \(\sin(u/2)\) and \(\sin((u-a)/2)\) have no common zero.  Their squared distance from their zero lattices has a positive minimum over one period, and their denominators are comparable to \((1+|u|)^4\); hence
\[
c(1+|u|)^{-4}\le g_{a,v}^{\sharp}(u)
\le C(1+|u|)^{-4}. \tag{58}
\]
Thus the second moment is finite.  Choose two different parameter pairs in (57), center both densities, and apply a common sufficiently small scale factor.  The resulting fixed densities \(g_1,g_2\) remain asymmetric and nonidentical; their characteristic functions are supported in one fixed interval \([-T_*,T_*]\), their second moments are at most \(M_e\), and the second density is centered.  The two-sided tail comparison (58) is preserved up to fixed constants.

Take \(p_x=|\mathcal X|^{-1}\) and \(f_x=f^\circ\) for every \(x\).  The declared restrictions on \(p_{\min}\), \(R_f\), and \(\kappa\) make these choices satisfy stratum positivity, the density Sobolev bound, compact support, and overlap.  Define
\[
m_{0,k}(a)=\frac12,
\qquad
m_{1,k}(a)=\frac12+\gamma\frac{h_k(a)}{f^\circ(a)},
\qquad
q_{j,k}=m_{j,k}f^\circ. \tag{59}
\]
Choose fixed \(\gamma>0\) smaller than \(R_q-\|f^\circ\|_{H^\beta}/2\).  The support of \(h_k\) is contained in the interior interval \(J\), where \(f^\circ\) has a positive minimum.  Since \(\beta>1/2\), (56) implies \(\|h_k\|_\infty\to0\), so, after reducing \(\gamma\) once, \(m_{1,k}\in[1/4,3/4]\) for all sufficiently large \(k\).  Equations (56) and (59) also give \(\|q_{j,k}\|_{H^\beta}\le R_q\).

Let \(V\) be uniform on \([0,1]\), independently of \((A,X,U_1,U_2)\), and define
\[
Y_j(a)=\mathbf 1\{V\le m_{j,k}(a)\},
\qquad Y_j=Y_j(A).
\]
This verifies consistency, bounded outcomes because \(B_Y\ge1\), conditional unconfoundedness, conditional Bernoulli marks, and causal-mean continuity.  Together with the preceding checks it proves \(P_{0,k},P_{1,k}\in\mathcal M\).  Their ADRFs satisfy
\[
\|\mu_{1,k}-\mu_{0,k}\|_{L^\infty(I)}
=\gamma\|h_k/f^\circ\|_{L^\infty(I)}
\ge ck^{1/2-\beta}. \tag{60}
\]

It remains to control the observed divergence.  Put
\[
d(s_1,s_2)=\int f^\circ(a)g_1(s_1-a)g_2(s_2-a)\,da,
\qquad
H_k(s_1,s_2)=\int h_k(a)g_1(s_1-a)g_2(s_2-a)\,da. \tag{61}
\]
Compact support of \(f^\circ\) and (58) imply
\[
d(s_1,s_2)
\ge c(1+|s_1|)^{-4}(1+|s_2|)^{-4}. \tag{62}
\]
The two-dimensional Fourier transform of \(H_k\) is
\[
\widehat H_k(s,t)=\widehat h_k(s+t)\varphi_1(s)\varphi_2(t), \tag{63}
\]
and it vanishes outside \([-T_*,T_*]^2\).  Differentiating \(\widehat h_k\), subtracting from \(e^{iua}\) its Taylor polynomial of order \(2k-j-1\), and using (50) gives, for \(0\le j\le4\) and \(|u|\le2T_*\),
\[
|\widehat h_k^{(j)}(u)|
\le C\|h_k\|_1\frac{C_0^{2k-j}}{(2k-j)!}
\le C\frac{C_0^{2k}}{(2k-4)!}. \tag{64}
\]
The concrete error densities in (58) have square-integrable products \(u^jg_\ell(u)\) for \(0\le j\le2\).  Thus their characteristic functions have the square-integrable weak derivatives needed below.  Apply Plancherel to (63), and to every mixed derivative through order two in each coordinate.  Leibniz' rule and (64) then give
\[
\int H_k(s_1,s_2)^2(1+|s_1|)^4(1+|s_2|)^4\,ds_1\,ds_2
\le C\frac{C_0^{4k}}{\{(2k-4)!\}^2}. \tag{65}
\]
Under the baseline law, conditional on any stratum, the observed densities for marks \(Y=1\) and \(Y=0\) are both \(d/2\).  Under the alternative their changes are respectively \(\gamma H_k\) and \(-\gamma H_k\).  Summing the two mark contributions and the common finite-stratum probabilities, then using (62)--(65), yields
\[
\chi^2(P_{1,k,\mathrm{obs}},P_{0,k,\mathrm{obs}})
=4\gamma^2\int\frac{H_k^2}{d}
\le C\frac{C_0^{4k}}{\{(2k-4)!\}^2}. \tag{66}
\]
Finally, the law of \((X,S_1,S_2)\) is unchanged because only the conditional mark was perturbed.  Hence all probability channels and the complete unmarked observed laws coincide.  Equations (60) and (66) prove every assertion.
